\scanpage{203}
\margin{202}
\hfill -- \mspage{56} -- \hfill IX

\uline{Lemme} 5.2.5. \uline{Soit} $u:G\longrightarrow H$ \uline{un homomorphisme de pro-$\ell$-groupes de Barsotti-Tate sur} $S$. \uline{Si} $u_o:G_o\longrightarrow H_o$ \uline{est nul}, \uline{alors} $u$ \uline{est nul}. \uline{Si} $u_o:G_o\longrightarrow H_o$ \uline{est une isogénie} (resp. \uline{un isomorphisme}), $u$ \uline{est une isogénie} (resp. \uline{un isomorphisme}).

Nous appliquerons en effet la première assertion de 5.2.5 au morphisme composé
\[
T_\ell(A^{\unsure{0|o}})^t\longrightarrow T_\ell(A^{\unsure{0|o}})^f\longrightarrow D(T_\ell(A'^{\unsure{0|o}})^f)\quad ,
\]
où le deuxième morphisme est celui déduit de l'accouplement (5.2.4.2), on trouve que ce composé est nul, donc que l'accouplement envisagé s'annule sur la partie torique à gauche. Par symétrie, il s'annule donc sur la partie torique à droite, donc il provient d'un accouplement sur les quotients, ou ``morceaux abéliens''
\[
T_\ell(A^{\unsure{0|o}})^{ab}\times T_\ell(A'^{\unsure{0|o}})^{ab}\longrightarrow T_\ell(\mathbf G_m)\quad ,
\]
lesquels quotients sont encore des pro-groupes de Barsotti-Tate, comme on vérifie aisément. Il reste à prouver que le morphisme correspondant à cet accouplement
\[
T_\ell(A^{\unsure{0|o}})^{ab}\longrightarrow D(T_\ell(A'^{\unsure{0|o}})^{ab})
\]
est une isogénie, ce qui résulte encore de 5.2.5.

Nous laisserons au lecteur la démonstration (facile) de 5.2.5, qui est un résultat valable sur toute base noethérienne connexe. Nous trouverons d'ailleurs plus bas (7.4.7), par une analyse un peu plus précise des biextensions $W_n$ des $(A_n^{\unsure{0|o}},A_n'^{\unsure{0|o}})$, une démonstration directe assez évidente des deux cas particuliers où nous avons eu à utiliser 5.2.5.
\margin{374}
\margin{374}
